\documentclass[11pt,a4paper]{article}

\usepackage[margin=2.3cm]{geometry}
\usepackage[T1]{fontenc}
\usepackage{tgpagella}
\usepackage[scale=0.92]{tgheros}
\usepackage{amsmath}
\usepackage{graphicx}
\usepackage{booktabs}
\usepackage{tabularx}
\usepackage{xltabular}
\usepackage{array}
\usepackage{float}
\usepackage[section]{placeins}
\renewcommand{\topfraction}{0.9}
\renewcommand{\bottomfraction}{0.8}
\renewcommand{\textfraction}{0.07}
\renewcommand{\floatpagefraction}{0.75}
\usepackage[font=small,labelfont=bf,skip=6pt]{caption}
\usepackage{xcolor}
\usepackage{enumitem}
\usepackage{titlesec}
\usepackage{fancyhdr}
\usepackage[hidelinks]{hyperref}

\definecolor{accent}{HTML}{0E7490}
\definecolor{ok}{HTML}{15803D}
\definecolor{part}{HTML}{B45309}
\definecolor{plus}{HTML}{1D4ED8}

\titleformat{\section}{\Large\sffamily\bfseries\color{accent}}{\thesection}{0.8em}{}
\titleformat{\subsection}{\large\sffamily\bfseries}{\thesubsection}{0.7em}{}
\titlespacing*{\section}{0pt}{18pt}{8pt}
\titlespacing*{\subsection}{0pt}{12pt}{5pt}

\setlist{itemsep=2pt,topsep=4pt}
\setlength{\parskip}{4pt}
\renewcommand{\arraystretch}{1.25}
\newcolumntype{L}{>{\raggedright\arraybackslash}X}

\newcommand{\done}{\textcolor{ok}{\textbf{Delivered}}}
\newcommand{\partly}{\textcolor{part}{\textbf{Partly}}}
\newcommand{\shot}[3][\textwidth]{%
  \begin{figure}[!htbp]\centering
  \includegraphics[width=#1]{figures/#2}
  \caption{#3}\label{fig:#2}
  \end{figure}}

\pagestyle{fancy}
\fancyhf{}
\fancyhead[L]{\small\sffamily ReactorGenome}
\fancyhead[R]{\small\sffamily CHE331 Innovation Project}
\fancyfoot[C]{\small\thepage}
\renewcommand{\headrulewidth}{0.4pt}

\begin{document}

\pagenumbering{roman}
\hypersetup{pageanchor=false}
\begin{titlepage}
  \centering
  \vspace*{3cm}
  {\sffamily\bfseries\fontsize{34}{40}\selectfont\color{accent} ReactorGenome\par}
  \vspace{0.4cm}
  {\Large\sffamily One Digital Twin for the Batch, CSTR, PFR and RTD Labs\par}
  \vspace{1.2cm}
  {\large CHE331 Innovation Project: Final Report\par}
  \vspace{2.5cm}
  {\large Mayur Dilip Pachpute\par}
  {Roll No. 240643\par}
  \vspace{0.6cm}
  {November 2026\par}
  \vfill
  \begin{minipage}{0.85\textwidth}
  \small
  \textbf{Summary.} ReactorGenome is a web application that sits beside the CHE331 reactor rigs. Students type the readings
  they already take; the app fits them to the reactor equations, shows what each textbook assumption costs, and lets the
  fitted kinetics drive reactor networks that cannot be built in the lab. All five parts of the proposal were built. Three
  points are only partly met and are listed openly in Section~\ref{sec:promised}. On the 2026 lab reports the app
  reproduces every reported rate constant and activation energy to four significant figures, and 134 automated tests
  check the engine, the pages and the exports.
  \end{minipage}
  \vspace{1.5cm}

  {\small Source code: \url{https://github.com/Mayurdp2403/ReactorGenome}\par}
\end{titlepage}
\hypersetup{pageanchor=true}

\tableofcontents
\newpage
\pagenumbering{arabic}

%=====================================================================
\section{The problem}

All four reactor experiments in CHE331 (Batch, CSTR, PFR and RTD) are run at fixed conditions: one reactor, three flow
rates, three temperatures and equal dilute feeds of NaOH and ethyl acetate. Readings are written down and analysed
later, so nobody sees what is happening inside the reactor while the rig is running. Questions such as why conversion
falls with flow, what a dead zone does, or why one residence time is not enough to reach steady state are answered in
words, never tested. In our own CSTR run the conductivity was still drifting when the ``steady'' readings were taken, and
28\,\textdegree C gave more conversion than 40\,\textdegree C. We had no way to find out why.

The proposal set out to build one web app, used on a laptop next to the rig with no new hardware, that turns the
manual readings into proper fits, live plots, assumption checks and what-if studies.

%=====================================================================
\section{What was built}

ReactorGenome is a Python application (Streamlit front end, NumPy/SciPy engine) with one page per experiment and a
Design Studio for reactor networks. It runs on a laptop with \texttt{run.bat}, and other laptops on the same network can
open the same address.

Every experiment page follows the same layout (Figure~\ref{fig:cstr_overview}): assumptions on the left, an animated
reactor in the middle, results and diagnostics on the right, the readings table below, then the analysis tabs. Every
control has a short help text.

\shot{cstr_overview}{CSTR page. Left: assumption switches and run conditions. Middle: animated tank with the simulated
start-up. Right: results, diagnostics and the kinetics fitted from the readings typed below (Ea~=~25.07~kJ/mol).}

The code has two layers:
\begin{itemize}
  \item \textbf{Engine} (\texttt{reactorgenome/engine}): the models and solvers, with no user-interface code. It covers
  units and bounds, rate laws, the stirred tank (steady, transient, non-ideal mixing and energy balance), batch fitting,
  plug flow with dispersion and laminar flow, RTD moments and model fits, the lab-sheet analysis and a flowsheet
  solver.
  \item \textbf{Interface} (\texttt{reactorgenome/ui}): one module per page, plus shared widgets, charts, cached engine
  calls and exports. The animation player and the drag-and-drop design board are small HTML/JavaScript components.
\end{itemize}

%=====================================================================
\section{Proposal against delivery}\label{sec:promised}

Table~\ref{tab:promised} takes each promise in the proposal in turn. ``Partly'' means the feature exists but does not
yet do everything the proposal describes.

{\small
\begin{xltabular}{\textwidth}{@{}>{\raggedright\arraybackslash}p{4.4cm} L >{\raggedright\arraybackslash}p{1.6cm}@{}}
\caption{What the proposal promised and what was delivered.}\label{tab:promised}\\
\toprule
\textbf{Proposal} & \textbf{Delivered} & \textbf{Status} \\
\midrule
\endfirsthead
\toprule
\textbf{Proposal} & \textbf{Delivered} & \textbf{Status} \\
\midrule
\endhead
\bottomrule
\endlastfoot
\multicolumn{3}{@{}l}{\textit{1. Enter the readings}}\\
Pick the rig and type readings at the times a stopwatch gives &
Batch, CSTR, PFR and RTD pages. The readings tables follow the lab sheets and accept CSV import. A stopwatch on
the CSTR page gives the reading schedule (2--5 settling times). & \done \\
\midrule
\multicolumn{3}{@{}l}{\textit{2. Fit the data properly}}\\
Treat CSTR readings as a transient and extrapolate to the true steady state &
Plateau from the simulated start-up shape, a single exponential, or the manual average. When reading times are typed,
each run gets a steady or not-steady verdict. & \done \\
Conductivity to conversion from ion concentrations, not a fixed $\Lambda_\infty$ &
Default $\Lambda_\infty = 0.366\,\Lambda_0$ from limiting ionic conductivities. A typed value is also accepted, and
$\Lambda_\infty$ can be estimated from the readings with a 95\% interval. & \done \\
Fit dead volume, bypass and mixing zones &
The bypass fraction is fitted from the readings. Dead volume, stagnant zones and tanks in series are set by the user and
compared against the readings. The RTD page measures $N$ and $Pe$ directly. & \partly \\
An error bar on every number &
95\% intervals on $k$, $E_a$ and $A$, error bars on measured conversions, and parity and residual plots. & \done \\
\midrule
\multicolumn{3}{@{}l}{\textit{3. Watch it live}}\\
Design plot, Arrhenius plot and transient redraw on each reading &
All plots and the fitted kinetics update as each cell is typed. & \done \\
Show the expected reading first, then judge ideal or not &
The app shows the outlet conductivity and conversion it expects for the current settings. Diagnostics after the fit
explain non-ideal behaviour. There is no step-by-step prediction before each individual reading. & \partly \\
\midrule
\multicolumn{3}{@{}l}{\textit{4. Change the assumptions}}\\
Switches for equimolar feed, irreversibility, ideal mixing, isothermal operation and steady state after one $\tau$ &
Each assumption is a switch on its page. The \emph{Assumption effects} tab shows the conversion with each one
dropped on its own (Figure~\ref{fig:cstr_assumptions}). & \done \\
\midrule
\multicolumn{3}{@{}l}{\textit{5. Try what the lab cannot do}}\\
Reactors in series in either order, unequal sizes, recycle, non-equimolar feeds, packed bed as a reactor &
Design Studio: a drag-and-connect board solved as one network, with presets and a mass-balance check. & \done \\
Non-isothermal scale-up with runaway and multiple steady states; Levenspiel plots and phase portraits &
Energy-balance tab with the heat-generation and heat-removal curves, stability of each steady state and phase
portraits. Levenspiel plots are on the PFR page and in the Design Studio. & \done \\
Batch $k$ improves the CSTR and PFR predictions in the same session &
Fitted second-order kinetics are published to every page and to the Design Studio. & \done \\
All four experiments on one conversion--$Da$ plot &
The CSTR page plots its measured runs on the $X$--$Da$ curve. The Design Studio overlays $k(T)$ from every rig and
shows the ideal $Da$ curves, but without the measured points from each rig. & \partly \\
\midrule
\multicolumn{3}{@{}l}{\textit{Deliverables}}\\
Four-tab app usable on a laptop or phone within the lab slot &
Four experiment pages plus the Design Studio; runs locally and over the lab network. Laptop use was tested; phone
layout was not. & \done \\
One-click export of the observation table, results, graphs and a diagnosis &
Each page exports a ZIP: an offline HTML report, CSV tables, every plot and every setting. The whole session saves as
a project file. & \done \\
\end{xltabular}}

%=====================================================================
\section{The application, page by page}

The screenshots below are the app running on the readings from the 2026 CHE331 lab reports. Those readings were
used only to check the mathematics. They are not built into the app, which starts with empty tables.

\FloatBarrier
\subsection{CSTR}

Readings are typed exactly as on the lab sheet (Figure~\ref{fig:cstr_readings}): one row per run with temperature,
total flow, $\tau$ and up to four outlet conductivities, plus a small table of $\Lambda_0$ for each temperature.

\shot{cstr_readings}{Readings table, laid out like the lab sheet's observation table.}

The \emph{Results} tab (Figure~\ref{fig:cstr_results}) collects everything the lab report needs on one screen:
$E_a$, $A$ and $k$ at each temperature with 95\% intervals, the design plot for each temperature and for all of them
together, the Arrhenius plot, conversion against flow and temperature, single-run $k$ against $\tau$, a parity plot
and residuals. The rate law used for the fit (order 1 to 3) is a setting.

\shot[0.86\textwidth]{cstr_results}{Results tab: kinetics with intervals and all of the report's plots on one screen.}

Below the plots, the fitted $k(T)$ predicts the conversion at any temperature, feed concentration and $\tau$
(Figure~\ref{fig:cstr_predictions}). The right-hand plot shows why one concentration cannot pin down the reaction
order: every order fits the runs at the $C_{A0}$ that was used, and the orders separate only away from it.

\shot{cstr_predictions}{What the fitted kinetics predict: $X$ against $\tau$, $k$ against $T$, and $X$ against $C_{A0}$
for orders 1 to 3.}

\shot{cstr_assumptions}{Assumption effects: conversion with each textbook assumption dropped on its own.}

The energy balance turns the tank into the scale-up case from the proposal. With a concentrated feed and a weak jacket
(Figure~\ref{fig:cstr_energy}) the heat-generation curve crosses the heat-removal line three times. The steady states at
31\,\textdegree C and 134\,\textdegree C are stable; the one at 80\,\textdegree C is not. The phase portrait shows which
start-up conditions run away to the hot state.

\shot{cstr_energy}{Energy balance for a 5~mol/L feed with a weak jacket ($UA = 40$~J/min/K, $E_a = 90$~kJ/mol,
$\Delta H = -100$~kJ/mol): three steady states and the phase portrait.}

\FloatBarrier
\subsection{PFR}

The PFR page treats the helical coil as a plug-flow reactor and lets the user switch on axial dispersion, laminar flow,
recycle, an adiabatic temperature rise and reaction that continues in the sample beaker. Coil Reynolds and Dean numbers
are worked out from the coil geometry.

\shot{pfr_overview}{PFR page with the kinetics fitted from the readings ($E_a = 12.70$~kJ/mol, as in the lab report).}

\shot[0.86\textwidth]{pfr_results}{PFR results: design plots, Arrhenius plot, conversion against flow and temperature,
single-run $k$, parity and residuals.}

\shot{pfr_flow_models}{The same $k$ under each flow picture, with the measured points at 50\,\textdegree C.}

\FloatBarrier
\subsection{Batch}

The batch page fits conductivity against time. The rate law is a choice, not an assumption: first order, second order,
reversible and a fitted power law are compared by AIC, with a clock offset for the seconds between mixing and starting
the stopwatch.

\shot{batch_overview}{Batch page at 50.2\,\textdegree C: fitted $k = 102.9$~L/mol/min with its interval.}

\shot{batch_fit}{Batch fit: integral and differential plots, residuals, and the rate laws ranked by AIC. The
differential method gives $n \approx 1.0$, the same as the lab report.}

\FloatBarrier
\subsection{RTD}

Up to four tracer runs share one time column. For each run the page gives the mean time, the variance, $N$ (both the
exact value and the nearest whole number), the Peclet number under three conventions, and fitted tanks-in-series and
dispersion models. It also warns when the curve was cut off before the tracer had left.

\shot{rtd_overview}{RTD page, stirred tank at 15~LPH. The diagnostic explains that $N = 1.5$ comes mostly from the cut-off tail.}

\shot{rtd_all_runs}{All four runs: tracer curves, $E(t)$ for each vessel, $E(\theta)$ and $F(t)$.}

\shot{rtd_model_fits}{Model fits ranked by error, with the residuals of the best model.}

\FloatBarrier
\subsection{Design Studio}

Blocks (feeds, CSTRs, PFRs, packed beds, splitters, products) are dragged onto a board, connected and solved as one
network. Recycle loops converge by Wegstein acceleration, and the mass balance closes to better than $10^{-8}$. The
same kinetics can come from any lab page.

\shot{design_studio}{Design Studio with a CSTR followed by a PFR: overall conversion 0.545, stream table and unit table.}

\FloatBarrier
\subsection{Export}

\shot{export}{Export tab: one ZIP with an offline HTML report, CSV tables, every plot and every setting used.}

%=====================================================================
\section{Validation}

The engine was checked in two ways: against closed-form results, and against the numbers in the 2026 lab reports
analysed with each report's own conventions. The full record, which a script regenerates from fresh calculations, is
in \nolinkurl{docs/VALIDATION.md}.

\begin{table}[H]
\centering\small
\caption{Engine against closed-form results (selected).}\label{tab:analytic}
\begin{tabular}{@{}lccc@{}}
\toprule
Case & Engine & Exact & Error \\
\midrule
Ideal CSTR, $Da = 1$ & 0.3820 & 0.3820 & $1\times10^{-16}$ \\
Ideal CSTR, $Da = 1$, $M = 2$ & 0.5858 & 0.5858 & $4\times10^{-14}$ \\
Plug flow, $X = Da/(1+Da)$ & 0.2308 & 0.2308 & $8\times10^{-12}$ \\
Dispersion, $Pe = 1000$ (plug-flow limit) & 0.2307 & 0.2308 & $9\times10^{-5}$ \\
Transient CSTR settles on the algebraic steady state & 0.4637 & 0.4637 & $5\times10^{-14}$ \\
Flowsheet CSTR $\rightarrow$ PFR against sequential solution & 0.5450 & 0.5450 & $5\times10^{-10}$ \\
RTD of 3 tanks: $N$ from moments & 3.000 & 3 & -- \\
\bottomrule
\end{tabular}
\end{table}

\begin{table}[H]
\centering\small
\caption{App against the 2026 lab reports, using each report's own conventions.}\label{tab:reports}
\begin{tabularx}{\textwidth}{@{}lLcc@{}}
\toprule
Experiment & Quantity & Report & App \\
\midrule
CSTR & $k$ at 31 / 46 / 60\,\textdegree C (L/mol/min) & 3.1629 / 5.8541 / 7.4362 & 3.1629 / 5.8541 / 7.4362 \\
     & $E_a$ (kJ/mol), $A$ (L/mol/min) & 25.07, $6.71\times10^4$ & 25.07, $6.71\times10^4$ \\
PFR  & $k$ at 40 / 50 / 60\,\textdegree C (L/mol/min) & 11.738 / 13.749 / 15.730 & 11.738 / 13.749 / 15.730 \\
     & $E_a$ (kJ/mol) & 12.70 & 12.70 \\
RTD  & CSTR 15 LPH: $t_m$ (min), $\sigma^2$ (min$^2$), $N$ & 8.244, 45.276, 1.501 & 8.244, 45.276, 1.501 \\
     & Packed bed 15 LPH: $Pe$, $\tau$ (min) & 35.93, 4.896 & 35.93, 4.896 \\
Batch & Reaction order (differential method) & 0.987 & 1.01 \\
\bottomrule
\end{tabularx}
\end{table}

The app reproduces every reported value. Applying the app's own methods to the same readings showed things the
reports could not:
\begin{itemize}
  \item \textbf{$\Lambda_\infty$ is the largest source of error.} Fitting one $\Lambda_\infty$ to all nine CSTR runs gives
  $\Lambda_\infty \approx 0.30\,\Lambda_0$, close to the ion estimate of $0.366\,\Lambda_0$ and far from the
  0.1~mS/cm used in the report. With it, $k$ at 31\,\textdegree C is about 8.5~L/mol/min rather than 3.16, and the
  single-run $k$ no longer drifts with $\tau$.
  \item \textbf{CSTR steady state.} Extrapolating each run to its plateau removes most of the
  ``28\,\textdegree C beats 40\,\textdegree C'' inversion described in the proposal. The inversion came from averaging
  readings taken before the tank had settled, not from the kinetics.
  \item \textbf{RTD tail.} The stirred-tank curves stop at 4\% of the peak. Adding the missing tail brings $N$ from 1.5 to
  about 1.2, close to one ideal tank.
  \item \textbf{Batch order.} After the mixing period the batch trace fits first order better than second order, which
  agrees with the report's own differential method.
\end{itemize}

%=====================================================================
\section{Beyond the proposal}

\begin{itemize}
  \item \textbf{Reaction order as a setting} (1 to 3) on the CSTR and PFR pages, with a table comparing how well each
  order fits and a plot showing where the orders separate.
  \item \textbf{Three ways to fit $k$} (conversion function on $\tau$, $\tau$ on conversion function, a free line), each
  with intervals, parity plots and residuals.
  \item \textbf{Single-run $k$ against $\tau$} as a consistency check: a sloping line points to $\Lambda_\infty$, a sampling
  delay, or runs that were not steady.
  \item \textbf{RTD extras:} tail correction with the tail's share of the tracer, three Peclet conventions, $\tau$ from
  $Pe$, $D_L$ from the bed geometry, and the conversion each measured vessel would give.
  \item \textbf{Sample delay in the PFR,} modelled explicitly instead of hidden in a free intercept of 7 to 16 minutes.
  \item \textbf{Input checks:} each page states the inputs it used and flags rows that cannot be analysed, such as a
  missing $\Lambda_0$ or readings beyond $\Lambda_\infty$.
  \item \textbf{Project files:} the whole session (tables, settings and the design board) saves to one file and loads back.
\end{itemize}

%=====================================================================
\section{Testing}

The test suite (\texttt{python -m pytest tests}, 134 tests) covers:
\begin{itemize}
  \item closed-form limits for every reactor model, unit conversions and bounds, invalid inputs and solver failures;
  \item flowsheets with recycle, the RTD moments and model fits;
  \item each lab report's numbers, reproduced to four significant figures;
  \item every page rendered with and without readings, including malformed tables;
  \item the exports, and saving and restoring a project.
\end{itemize}
\texttt{tools/make\_validation.py} regenerates the validation record from fresh calculations.

%=====================================================================
\section{Limitations}

\begin{itemize}
  \item The three ``Partly'' items in Table~\ref{tab:promised}: dead volume and stagnant zones are compared rather than
  fitted, there is no prediction step before each individual reading, and the common $Da$ plot does not yet carry the
  measured points from every rig.
  \item The layout has not been tested on a phone.
  \item Only second-order kinetics are shared between pages. Fits with another order stay on the page that made them.
  \item The energy balance is checked against its isothermal limit and for multiple steady states, not against measured
  temperatures, because the rigs are run isothermally.
  \item The chemistry is NaOH with ethyl acetate throughout, as in the course.
\end{itemize}

%=====================================================================
\section{Running the app}

\begin{verbatim}
git clone https://github.com/Mayurdp2403/ReactorGenome.git
cd ReactorGenome
python -m pip install -r requirements.txt
python -m streamlit run app.py
\end{verbatim}
On Windows, \texttt{run.bat} does the same. The screenshots in this report are regenerated with
\texttt{python docs/report/make\_figures.py}.

\end{document}
